\section{Ablations}\label{sec:ablations}

Do the specific design choices inside S3-FIFO, SIEVE, and W-TinyLFU matter?
Each of these three policies packages a concrete mechanism: S3-FIFO's
small-queue admission ratio, SIEVE's visited-bit lazy promotion, and
W-TinyLFU's optional Doorkeeper Bloom filter. We run three ablation axes,
holding the policy fixed and varying one design parameter at a time across
$\alpha \in \{0.6, \ldots, 1.2\}$ at $\mathrm{cache\_frac}=0.01$ on both
workloads. The three ablations live together in this single section
deliberately, as a rigor signal: we're not cherry-picking one result, we're
reporting all three mechanism decompositions side-by-side.

\subsection*{S3-FIFO small-queue ratio}

Yang et al.\ (\cite{yang2023fifo}) recommend a $10\%$ small-queue default
in the SOSP'23 S3-FIFO design. Does this default generalize across
workloads? We run the three-queue policy at small-queue ratios
$\{5\%, 10\%, 20\%\}$ across the full $\alpha$ grid on both workloads.

\begin{figure}[h]
\centering
\includegraphics[width=\textwidth]{../results/congress/figures/ablation_s3fifo.pdf}
\caption{S3-FIFO small-queue ratio ablation --- $5\%$, $10\%$, and $20\%$
variants on Congress (left) and Court (right) across
$\alpha \in \{0.6, \ldots, 1.2\}$ at $\mathrm{cache\_frac}=0.01$. A
two-panel shared-$y$-axis view; smaller small-queue ratios win
monotonically.}
\label{fig:ablation_s3fifo}
\end{figure}

\textbf{Finding:} smaller small-queue ratios win monotonically across every
$\alpha$ on both workloads --- a $6.3$\,pp gap at Court $\alpha=1.2$
($5\%$ variant beats $20\%$ variant) versus $1.2$\,pp on Congress at the
same $\alpha$. Court's heavier tail amplifies the advantage of a smaller
probation pool. This is a publishable result against the $10\%$ paper
default: practitioners deploying on a heavy-tailed workload should reduce
the small-queue ratio toward $5\%$; practitioners on a near-uniform
workload can follow the paper default with little penalty.

\subsection*{SIEVE visited-bit (lazy promotion)}

Zhang et al.\ (\cite{zhang2024sieve}) claim lazy promotion --- scanning
forward from a pointer and clearing visited bits rather than eagerly
promoting on each hit --- is the central SIEVE innovation. Does it
quantitatively matter? We toggle the \texttt{promote\_on\_hit} flag and
re-sweep $\alpha$.

\begin{figure}[h]
\centering
\includegraphics[width=\textwidth]{../results/congress/figures/ablation_sieve.pdf}
\caption{SIEVE visited-bit ablation --- standard SIEVE (solid) vs.\
SIEVE-NoProm (dashed) across $\alpha$ at $\mathrm{cache\_frac}=0.01$ on
both workloads. SIEVE-NoProm monotonically loses at every point.}
\label{fig:ablation_sieve}
\end{figure}

\textbf{Finding:} SIEVE-NoProm monotonically loses to SIEVE at every
$\alpha$ on both workloads. The peak gap is $+15.4$\,pp on Congress at
$\alpha=1.0$ and $+11.0$\,pp on Court at $\alpha=1.1$. This is among the
largest ablation deltas in our study and empirically confirms the NSDI'24
paper's claim that lazy promotion is the dominant contributor to SIEVE's
scan-resistance --- the visited-bit scanner is what allows SIEVE to behave
differently from a plain FIFO despite its single-queue structure.

\subsection*{W-TinyLFU $\pm$ Doorkeeper}

Einziger and Friedman (\cite{einziger2017tinylfu}) describe an optional
pre-CMS Bloom filter called the ``Doorkeeper'' that is meant to absorb
one-hit-wonder traffic before it pollutes the Count-Min Sketch. Caffeine's
production W-TinyLFU (\cite{caffeine2024}) omits it. We compare W-TinyLFU
with and without Doorkeeper across $\alpha$ on both workloads.

\begin{figure}[h]
\centering
\includegraphics[width=\textwidth]{../results/congress/figures/ablation_doorkeeper.pdf}
\caption{W-TinyLFU $\pm$ Doorkeeper ablation across $\alpha$ at
$\mathrm{cache\_frac}=0.01$. Shared brown color; solid = baseline
W-TinyLFU, dashed = W-TinyLFU$+$DK variant. Congress panel shows
essentially overlapping curves; Court panel shows a small direction-swap
near the $\alpha$ extremes.}
\label{fig:ablation_dk}
\end{figure}

\textbf{Finding:} context-dependent with a small effect at the $\alpha$
extremes. Congress stays within $\pm 0.22$\,pp across every $\alpha$
(functionally noise). Court shows a peak $-0.72$\,pp win for the $+$DK
variant at $\alpha=1.1$, but a $+0.52$\,pp \emph{loss} at $\alpha=0.6$.
\textbf{Interpretation:} Doorkeeper is a marginal hedge against one-hit-
wonder spikes on a heavy-tailed workload, not a default. The
direction-swap at low $\alpha$ confirms the Caffeine production team's
decision to omit it --- a workload without a real frequency gradient has
nothing for Doorkeeper to absorb, and the filter's own false-positives
become a small penalty.

\subsection*{Ablation takeaway}

Reading the three ablations together: \textbf{SIEVE's lazy promotion is
load-bearing} (the largest ablation delta in our study, $+15.4$\,pp peak);
\textbf{S3-FIFO benefits from a smaller small-queue than the SOSP'23
default} (actionable advice for practitioners on heavy-tailed workloads);
and \textbf{Doorkeeper is context-dependent} (workload-sensitive hedge
rather than a default). Together these results justify the claim that our
policy ordering in \S\ref{sec:policies} is stable across reasonable
parameter variations within each policy --- the ordering is a property of
the policy families, not of the specific parameter choices we happened to
ship in each implementation.
